% Part: lambda-calculus
% Chapter: introduction
% Section: lambda-computable

\documentclass[../../../include/open-logic-section]{subfiles}

\begin{document}

\olfileid{lam}{int}{cmp}
\olsection{\usetoken{S}{lambda definable} ప్రమేయాలు గణనీయమైనవి}

\begin{thm}
\ollabel{thm:lambda-computable}
పాక్షిక ప్రమేయం~$f$ ఒక లాంబ్డా పదంచే \tetoken{లాంబ్డాతో నిర్వచించబడిన}{!!{lambda defined}}దైతే, అది
గణనీయమైనది.
\end{thm}

\begin{proof}
ప్రమేయం~$f$ లాంబ్డా పదం~$X$చే \tetoken{లాంబ్డాతో నిర్వచించబడిన}{!!{lambda defined}}దని అనుకుందాం.
$f$ను గణించే ఒక అనౌపచారిక ప్రక్రియను వర్ణిద్దాం. ఇన్‌పుట్‌గా $m_0$,
\dots,~$m_{n-1}$ ఇచ్చినప్పుడు, $X \num m_0 \ldots \num
m_{n-1}$ పదాన్ని రాయండి. మొదట అసలు పదానికి గల ఒక-దశ తగ్గింపులన్నిటినీ
రాసి ఒక వృక్షాన్ని నిర్మించండి; వాటి కింద ఆ పదాల ఒక-దశ తగ్గింపులన్నిటినీ
(అంటే అసలు పదపు రెండు-దశ తగ్గింపులను) రాయండి; అలాగే కొనసాగించండి.
ఎప్పుడైనా ఒక సంఖ్యాంకాన్ని చేరితే, దానినే సమాధానంగా తిరిగి ఇవ్వండి;
లేకపోతే ప్రమేయం అనిర్వచితం.

చర్చ్ సిద్ధాంతప్రతిపాదనను ఆశ్రయిస్తే ఈ ప్రమేయం గణనీయమని తెలుస్తుంది.
ఈ సిద్ధాంతాన్ని నిరూపించడానికి మెరుగైన మార్గం, ఈ అన్వేషణ ప్రక్రియకు
పునరావృత్త వర్ణన ఇవ్వడం. ఉదాహరణకు, ``$\fn{IsASubterm}$,''
``$\fn{Substitute}$,'' ``$\fn{ReducesToInOneStep}$,''
``$\fn{ReductionSequence}$,'' ``$\fn{Numeral}$,'' మొదలైన ఆదిమ
పునరావృత్త ప్రమేయాలు, సంబంధాల అనుక్రమాన్ని నిర్వచించవచ్చు. అప్పుడు
$f(m_0, \dots, m_{n-1})$ను గణించే పాక్షిక పునరావృత్త ప్రక్రియ,
$X \num{m_0} \dots \num{m_{n-1}}$తో మొదలై ఒక సంఖ్యాంకంతో ముగిసే
ఒక-దశ తగ్గింపుల అనుక్రమం కోసం అన్వేషించి, ఆ సంఖ్యాంకానికి అనుగుణమైన
సంఖ్యను తిరిగి ఇవ్వడం. వివరాలు దీర్ఘమైనవి, శ్రమతో కూడినవి; కానీ
మిగతా విషయాల్లో మామూలైనవే.
\end{proof}

\end{document}
